% Introduction for the EJW paper. Numbers from sections/numbers-lttv.tex;
% literature figures quoted from the sources with locators. Draft, 2026-10-03.

\section{Introduction}
\label{sec:intro}

On 12 April 2016 a committee of the House of Commons heard that the Canadian Radio-television and Telecommunications Commission's (CRTC's) new rules for television would, \enquote{by 2020}, cost \enquote{some 15,130 media jobs}, \enquote{all of this as the direct result of Let's Talk TV regulatory changes} \citep[time marks 0900--0905]{commons2016chpc8}. The number came from an economic-impact study by the consultancy Nordicity and Peter Miller, an engineer and communications lawyer, prepared for ACTRA, the Canadian Media Guild, the Directors Guild of Canada, Friends of Canadian Broadcasting and Unifor \citep{nordicity2015television}. The CRTC had told distributors to offer a \$25 entry-level package and to sell discretionary channels singly or in small bundles \citep{crtc2015worldofchoice}. The study modelled what that would do to channel revenue, programming spending, jobs and GDP through 2020.

Outcomes for those years are now public. I found no published check of the forecast, in the scholarly literature, in the trade press, or in the authors' later reports for the CRTC (the searches are listed in the replication files). This paper supplies one.

Economists have long compared what was predicted before a regulation with what happened after it. \textcite[ii]{harrington1999accuracy} found that ex ante estimates of direct costs exceeded actual costs for 12 of 25 environmental and occupational-safety rules and fell short for 6, with errors \enquote{driven by both baseline and compliance issues}. A later set of case studies found \enquote{a slight tendency to overestimate both costs and benefits} \citep[285]{morgenstern2018retrospective}. \textcite{simpson2014overestimate} warns that a preponderance of overestimates need not show bias, since skewed cost distributions could produce one anyway. His review notes that the few studies of industry's own estimates found them above outcomes in most cases \citep[319--320]{simpson2014overestimate}.

The work I've cited concerns mainly regulators' estimates of compliance costs and benefits. The forecasts that interested parties put to legislators are a different object. They predict economy-wide jobs and output through input-output multipliers, they travel through testimony rather than regulatory filings, and I know of no routine for revisiting them. \textcite[4]{manski2011certitude} observes that in policy analysis \enquote{Point predictions are common and expressions of uncertainty are rare}. A forecast delivered to a parliamentary committee as a single number of jobs that \enquote{will be lost} is a clear case.

The closest precedent is \textcite{kane2026sffa}, who checked the enrolment forecasts made by universities and by Harvard's expert during the litigation that ended in the US Supreme Court's decision in \textit{Students for Fair Admissions v.\ Harvard}. Two of his choices carry over. He reports results against more than one baseline, and he separates the expert's technical estimate, stated in units no public statistic matches, from the forecasts the institutions put to the Court.

This paper is built on the same separation. The report's technical claim is a conditional projection: if the reforms were implemented and its assumptions held, revenue, programming spending and jobs would follow given paths. The public argument is what the committee was told would happen. The two can fare differently, and here they do.

I fixed the forecast record before downloading any outcome data, ran a fake-data design analysis to see which quantities could tell a world with the forecast effect from one without it, and wrote down how each result would be read \citep{gelman2014types, gelman2020workflow}. The main outcome measure is a single number, \(k\): the observed shortfall below the report's own no-reform baseline, as a multiple of its forecast impact.

Specialty and pay channels' revenue, the link on which most of the job figure depends, didn't fall: \(k=\lttvKSpec\), with an interval from \lttvKSpecLo{} to \lttvKSpecHi{}. That verdict is fragile. It depends on how wrong the report's baseline is taken to be, and it becomes inconclusive under \lttvNInconclusiveSpec{} of the \lttvNCalibrations{} calibrations I examined. Distributors' revenue fell further than forecast, though with a wide interval (\(k=\lttvKBDU\), from \lttvKBDULo{} to \lttvKBDUHi{}).

The break in the chain doesn't depend on the baseline at all. Applied to distributors' actual revenue fall, the report's own pass-through rule implies a cut of about \$\lttvPassImplied{} million in what they paid Canadian channels; those payments rose by \$\lttvPayCanIncrease{} million. Uptake of the entry-level package, a key input, was \lttvUptakeJuneShare{} of subscribers in mid-2016 against an assumed \lttvByopSixteen{} for that year.

The testimony went beyond the report. It turned an economy-wide estimate, more than half of it spin-off employment in other sectors, into media jobs lost as the direct result of the policy, and it cited a share of programming spending that the report's own figures don't reproduce.

Section~\ref{sec:lttv-forecast} sets out the decision, the forecast and the testimony. Section~\ref{sec:lttv-method} describes the check, section~\ref{sec:lttv-results} the results, and section~\ref{sec:lttv-links} the links in the chain. Section~\ref{sec:discussion} discusses what the case shows and what it doesn't.
